% Standalone section fragment; requires amsmath, amssymb, amsthm.
\section{A generic obstruction for coordinate-sign networks}
\label{sec:sign-cube-obstruction}

Adding many full matchings does not by itself remove the diagonal
obstruction.  The following example has numerous cycles, but its relevant
Fourier component still has only two independent flow channels.

Let \(K\) be an imaginary quadratic field and let \(E\) be a complex
elliptic curve with \(\operatorname{End}^0(E)=K\).
Put \(m=2n\), let \(A=E^{m}\) have its product principal polarization
\(\theta\), and use the diagonal balanced action with eigenspaces
\(U=\langle e_1,\ldots,e_n\rangle\) and
\(V=\langle e_{n+1},\ldots,e_m\rangle\).  Its polarized tangent space is
\[
 T=\left\{\begin{pmatrix}0&B\\B^t&0\end{pmatrix}:B\in M_n(\mathbb C)\right\}.
\]
Let \(D_\sigma\) run through all permutations of the distinct eigenvalues
\(1,-1,\ldots,n,-n\).  For a real orthogonal frame \(g\), form the complete
permutation packets
\[
 D_{\epsilon,\sigma}=\epsilon gD_\sigma g^t\epsilon,
 \qquad
 \epsilon\in\Gamma:=\{\operatorname{diag}(\pm1,\ldots,\pm1)\}/\{\pm I\}.
\]
Retain all transposition edges inside each packet, and all full matchings
induced by the coordinate sign reflections \(r_l\).
The graph on the packets is the folded \(m\)-cube.  Its parity is
well defined because \(m\) is even, and combining this parity with
permutation parity gives a bipartition.

For an edge \(e=\{x,y\}\), write \(\ell_e\subset H^2(A,\mathcal O_A)\)
for its antiholomorphic surface wedge line.  The trace-incidence
condition for weights \(m_x\) and a tangent \(\kappa\) is the existence
of edge values \(\gamma_e\in\ell_e\) such that
\[
 m_x\,\kappa\mathbin{\lrcorner}D_x
   =\sum_{e\ni x}\varepsilon(x,e)\gamma_e,
\]
with opposite signs at the two endpoints.  These are necessary
conditions for an unfiltered deformation of an ordinary extension
supported on the graph.

\begin{theorem}\label{thm:sign-cube-obstruction}
Suppose \(n\ge3\).  There is a nonempty Zariski-open set of real
orthogonal frames \(g\), containing rational frames, for which the
following assertion holds for every positive integer \(N\) and every
choice of real weights \(m_x\).  If the weighted complete character is
\[
 \sum_x m_x e^{ND_x}=R F_n(N\theta)+cw,
 \qquad c\ne0,\quad 0\ne w\in W_K(A),
\]
then the trace-incidence condition fails for some \(\kappa\in T\).
In particular no ordinary extension with these multiplicities and
this support graph can deform in every polarized Weil direction,
even after its filtration is forgotten.
\end{theorem}

For real frames this is a linear-algebraic obstruction.  For rational
frames, choose \(N\) so that every \(ND_x\) is an integral Chern class;
positive integral weights then describe multiplicities of line bundles
in an ordinary extension.  The weights need not be positive for the
obstruction argument.  The theorem concerns this family of coordinate-sign networks,
and makes no assertion about arbitrary networks of full matchings.

\begin{proof}
Let \(\chi_{\rm all}(\epsilon)=\prod_l\epsilon_l\), a character of
\(\Gamma\).  In \(H^{p,p}\), the coefficient
\(z_I\wedge\overline z_J\) has sign character
\(\prod_{l\in I\mathbin{\triangle}J}\epsilon_l\).
Consequently the \(\chi_{\rm all}\)-component consists exactly of
the middle-degree coefficients
\[
 |I|=|J|=n,\qquad I\cap J=\varnothing,\qquad I\cup J=\{1,\ldots,m\}.
\]
They are primitive, since wedging with \(\theta\) repeats a factor.
The fixed Weil plane lies in this component.  Not every class in
this component is a Weil class.

The graph and \(T\) are stable under \(\Gamma\).  If the trace equations
are solvable for every tangent, averaging the weights and flows with
\(\chi_{\rm all}\), while transporting the tangent, preserves
solvability.  We may therefore consider weights
\(\chi_{\rm all}(\epsilon)f(\sigma)\).  Their primitive character must
be nonzero if the prescribed Weil term is nonzero.

Fix a basic tangent
\(B_{ij}=E_{ij}+E_{ji}\), taking \(i=1\) and \(j=n+1\).
The open set of frames below is chosen for this pair.
In the frame \(\epsilon g\), its matrix is
\[
 g^t\epsilon B_{ij}\epsilon g
   =\epsilon_i\epsilon_jg^tB_{ij}g.
\]
The resulting scalar Fourier mode is
\(\psi=\chi_{\rm all}\chi_{ij}\).
Along the coordinate-reflection matching at \(l\), the transported
bridge wedge changes sign at the other endpoint.  Its scalar
coefficient must therefore be invariant under \(r_l\).
In Fourier mode \(\psi\), this permits a channel only if
\(\psi(r_l)=1\), which holds exactly for \(l=i,j\).

Set \(u=g^te_i\), \(v=g^te_j\).  At an internal transposition
\(s=(kl)\), the tangent coefficient is
\(u_kv_l+v_ku_l\); the two permitted bridge coefficients are
\(u_ku_l\) and \(v_kv_l\).  On functions of the permutation vertex,
define \((R_s f)(\sigma)=f(s\sigma)\).  Add the \((k,l)\)-coordinate equations
at \(\sigma\) and \(s\sigma\).  Their unique internal edge cancels.
The nonzero eigenvalue difference can be divided out, giving
\begin{equation}
 (1-R_s)
 \left[(u_kv_l+v_ku_l)f-u_ku_l z_1-v_kv_l z_2\right]=0
 \quad\text{for every }s=(kl).
 \label{eq:sign-cube-reduction}
\end{equation}
where \(z_1,z_2\) are arbitrary scalar potentials on the permutation
vertices.  Thus this reduction has retained all possible reverse
flow coefficients.

These equations commute with the right permutation action, so they
may be examined separately in each left irreducible constituent.
Only the middle primitive constituent
\(\mathcal P_n\simeq S^{(n,n)}\) contributes a primitive middle
character.  In the square-free algebra generated by the \(x_l\),
this is the span of the products
\[
 p_M=\prod_{\{a,b\}\in M}(x_a-x_b)
\]
over perfect matchings \(M\) of \(\{1,\ldots,2n\}\).
It occurs in the diagonal degree-\(n\) module and in no other
diagonal degree.

We use the following elementary observation.  Suppose a representation
satisfies equations
\[
 (1-R_s)\left[a_sf-\sum_{r=1}^q c_{r,s}z_r\right]=0.
\]
For a perfect matching \(M\), its \(n\) transpositions commute.
Apply \(D_M=\prod_{s\in M}(1-R_s)\) to obtain
\begin{equation}
 a_sD_Mf-\sum_{r=1}^q c_{r,s}D_Mz_r=0,\qquad s\in M.
 \label{eq:sign-cube-matching}
\end{equation}
If the scalar coefficient matrix of these \(n\) equations has
column rank \(q+1\), then \(D_Mf=0\).
If this holds for every perfect matching, the
\(S^{(n,n)}\)-component of \(f\) is zero.  Indeed \(D_M\) is a
nonzero column antisymmetrizer on that irreducible module, and the
joint kernel of all its conjugates is therefore zero.  Equivalently,
the products \(p_M\) span the primitive module, and the operators
\(D_M\) are self-adjoint in the coordinate-subset inner product.

For \eqref{eq:sign-cube-reduction}, assume \(u_k\ne0\) and put \(t_k=v_k/u_k\).
Dividing the row associated with \(\{k,l\}\in M\) by \(u_ku_l\)
turns its scalar coefficients into
\begin{equation}
 \bigl[t_k+t_l,\ -1,\ -t_kt_l\bigr].
 \label{eq:sign-cube-rank}
\end{equation}
For \(n\ge3\), these rows have rank three for every perfect matching
for a generic tuple \(t\).  Each failure is an algebraic condition.
For each fixed matching it is proper: choose three of its disjoint
pairs so that their points
\((t_k+t_l,t_kt_l)\) are not collinear.
There are only finitely many matchings, so the simultaneous
nonvanishing conditions define a nonempty Zariski-open set of
tuples.

These conditions are compatible with \(u,v\) being real orthonormal
vectors.  The set of tuples \(t\) for which there are positive
weights \(w_k\) satisfying
\[
 \sum w_k=1,\qquad \sum w_kt_k=0,\qquad \sum w_kt_k^2=1
\]
contains a nonempty Euclidean open set.  For example, for six
coordinates take values
\(-2,-1,-\tfrac12,\tfrac12,1,2\), with respective symmetric
weights \(1/12,1/12,1/3,1/3,1/12,1/12\).
Additional coordinates may be assigned small positive weights,
with a small adjustment of the first six.  The point \((0,1)\)
is in the interior of the corresponding convex hull of
\((t_k,t_k^2)\), so feasibility persists under small perturbations.
This open set meets the complement of the finitely many algebraic
exceptional sets.  Setting \(u_k=\sqrt{w_k}\) and
\(v_k=u_kt_k\) gives orthonormal vectors, which can be completed
to an orthogonal frame.

It follows that the required rank conditions define a nonempty
Zariski-open subset of the orthogonal group.  Rational orthogonal
matrices are Zariski dense, for example by the rational Cayley
parametrization, so rational frames occur.  We may also exclude
unintended packet coincidences: for each noncentral coordinate
sign matrix, the condition that its conjugate by \(g\) is a signed
permutation matrix is a proper algebraic condition; there are
only finitely many such conditions.

For these frames, \eqref{eq:sign-cube-matching}--\eqref{eq:sign-cube-rank} kill the primitive
\(\chi_{\rm all}\)-component already for the chosen basic tangent.
This contradicts the required nonzero Weil character and proves
the theorem.
\end{proof}

